\documentclass[10pt,a4paper]{amsart}
\usepackage[margin=19mm]{geometry}
\usepackage{amsmath,amssymb,mathtools}
\usepackage[scr=rsfs,cal=euler]{mathalfa}
\usepackage{xfrac}
\usepackage[unicode,hidelinks,bookmarks=false]{hyperref}
\newcommand{\ie}{i.e.\ }
\newcommand{\cf}{cf.\ }
\newcommand{\resp}{respectively\ }
\newcommand{\Subsection}{\subsection}
\providecommand{\nobreakdash}{\nobreak\hbox{-}}
\setlength{\emergencystretch}{2em}
\multlinegap0pt
\usepackage{xcolor}
\usepackage{amssymb}
\newtheorem{prop}{Proposition}
\title{What do the Enneper surface and helix surfaces have in common?}
\author{Pascual Lucas, Jos\'e Antonio Ortega-Yag\"ues}
\date{Source: arXiv:2601.19253v1 (CC BY 4.0)}
\begin{document}
\maketitle

\noindent\emph{Source and licence.} What do the Enneper surface and helix surfaces have in common?. Version arXiv:2601.19253v1 (CC BY 4.0), distributed by arXiv under the Creative Commons Attribution 4.0 International licence, http://creativecommons.org/licenses/by/4.0/, as recorded in the arXiv licence metadata for this exact version and shown on the abstract page https://arxiv.org/abs/2601.19253v1. Full licence notice: \texttt{licenses/arxiv-2601.19253-CC-BY-4.0.txt}.\par\smallskip

\noindent\textit{Editorial scope.} Counted unit: the article's prop p1.1 (source0.tex lines 218--226). The article's complete original proof of it is delivered with it (source0.tex lines 227--234, one \texttt{proof} environment of 8 lines). No mathematics is written, repaired or re-derived: the statement and the proof are the article's own lines, in the article's own notation, and no retained line is substituted or re-flowed.  Retained uncounted as the source-local context the proof needs: lines 132--135.  Nothing was omitted from any retained span, so the package carries no omission record.  No retained line contains a citation, so the wrapper carries no bibliography and no bibliography entry.  No retained line contains a figure environment, an image-inclusion command or drawing code, so no image is delivered and the package declares no asset dependency.  Editorial formatting only, in addition: an AMS \texttt{amsart} preamble that restates the article's own declarations and macros used by the retained lines, namely the environments \texttt{prop} on separate counters, together with the standard packages those lines need.  The retained environments print in this document as Proposition 1 in order of appearance, whereas the article prints the same items with the numbering of its own full document.  The complete native source of the article as distributed in the arXiv e-print for this exact version is bundled unchanged as \texttt{sources/193445/source0.tex}.
\begin{equation}\label{kntg}
\kappa_n=\kappa_1\cos^2\phi+\kappa_2\sin^2\phi,\qquad
\tau_g=(\kappa_1-\kappa_2)\cos\phi\sin\phi.
\end{equation}

\begin{prop}\label{p1.1}%
Let $\gamma$ be an arclength parametrized curve in a surface $M$. If $\gamma$ is an isogonal line (which is not a line of curvature), then the following conditions are equivalent:\vspace*{-\topsep}
\begin{enumerate}\itemsep0pt\def\labelenumi{\roman{enumi})}
\item[C1)]
$\kappa_1$ and $\kappa_2$ are linearly dependent functions along $\gamma$.
\item[C2)]
$\tau_g$ and $\kappa_n$ are linearly dependent functions along $\gamma$.
\end{enumerate}
\end{prop}

\begin{proof}
From equations (\ref{kntg}), it is easy to see that, for any constants $a,b,c,d$ and $\phi$, we have
\begin{align}
\sin\phi\cos\phi(a+b)\kappa_n+(a\sin^2\phi-b\cos^2\phi)\tau_g &= \sin\phi\cos\phi(a\kappa_1+b\kappa_2).\label{(1)}\\
(d\cos\phi+c\sin\phi)\cos\phi\kappa_1+(d\sin\phi-c\cos\phi)\sin\phi\kappa_2 &= c\tau_g+d\kappa_n.\label{(2)}
\end{align}
From these equations, it is not difficult to see that \emph{C1} and \emph{C2} are equivalent conditions.
\end{proof}

\end{document}
